% =============================================================
% 4.4 การออกแบบเชิงการศึกษาและการประเมินผล
% โครงงาน: การออกแบบและพัฒนาเกมเดินเรือด้วยระบบสั่งการด้วยเสียงเพื่อส่งเสริมทักษะการสื่อสาร
% =============================================================

\section{การออกแบบเชิงการศึกษาและการประเมินผล}

เนื่องจากโครงงานนี้เป็นเกมเพื่อการเรียนรู้แบบจริงจัง (Serious Game) การออกแบบจึงบูรณาการกลไกของเกม (Gamification) เข้ากับทฤษฎีการเรียนรู้ภาษา เพื่อให้ผู้เล่นได้รับการพัฒนาทักษะเชิงปฏิบัติอย่างมีประสิทธิภาพ\cite{sadigzade2025} โดยแบ่งออกเป็น 2 ส่วน ดังนี้

\subsection{การออกแบบเชิงการศึกษา (Educational Design)}

การออกแบบกลไกเพื่อส่งเสริมการเรียนรู้ประกอบด้วยแนวคิดหลัก ดังนี้

\begin{enumerate}[label=\arabic*)]
    \item \textbf{การเรียนรู้ตามสถานการณ์ (Situated Learning):} ผู้เล่นจะไม่ได้เรียนรู้คำศัพท์จากการท่องจำ แต่จะได้เรียนรู้ผ่านการจำลองสถานการณ์เดินเรือจริง (Maritime Communication) เช่น การสั่งหันหางเสือหรือการปรับความเร็ว เพื่อให้เกิดความเข้าใจในบริบทและสามารถนำไปประยุกต์ใช้ได้ทันที\cite{sadigzade2025}
    \item \textbf{การลดตัวกรองทางอารมณ์ (Affective Filter Hypothesis):} การให้ผู้เล่นฝึกเปล่งเสียงสั่งการกับ "ลูกเรือปัญญาประดิษฐ์ (AI Crew)" แทนการสนทนากับมนุษย์จริง จะช่วยสร้างพื้นที่ปลอดภัย (Safe Space) ลดความกลัวความผิดพลาด และบรรเทาความวิตกกังวลในการใช้ภาษาต่างประเทศ (Language Anxiety)\cite{sadigzade2025}
    \item \textbf{การเรียนรู้จากข้อผิดพลาด (Learning from Failure):} เมื่อผู้เล่นเปล่งเสียงไม่ชัดเจนหรือผิดหลักไวยากรณ์ ระบบจะไม่ลงโทษด้วยการตัดจบทันที แต่ลูกเรือ AI จะใช้กลไกการทวนถาม (Clarification Request) ซึ่งเป็นการให้ข้อมูลป้อนกลับ (Feedback) เพื่อกระตุ้นให้ผู้เล่นแก้ไขการพูดด้วยตนเอง\cite{zargham2024}
\end{enumerate}

\subsection{แนวทางการประเมินผล (Evaluation Approach)}

การประเมินผลของเกมถูกออกแบบให้ครอบคลุมทั้งประสิทธิภาพของเทคโนโลยีและผลสัมฤทธิ์ด้านการเรียนรู้ของผู้เล่น ดังนี้

\begin{itemize}
    \item \textbf{การประเมินผลภายในเกม (In-game Evaluation):} ประเมินจากความแม่นยำในการออกเสียงคำสั่ง (Speech Accuracy) ความรวดเร็วในการสั่งการแก้ไขปัญหาเฉพาะหน้า และสภาพความสมบูรณ์ของตัวเรือ (Hull Integrity) ที่เหลืออยู่เมื่อจบภารกิจ
    \item \textbf{การประเมินประสิทธิภาพทางเทคนิค (Technical Evaluation):} ประเมินอัตราความแม่นยำในการแปลคำสั่งเสียงของระบบ Speech-to-Text และวัดค่าความหน่วง (Latency) ของระบบตั้งแต่ผู้เล่นพูดจบจนถึงตอนที่เรือเริ่มตอบสนอง
    \item \textbf{การประเมินประสบการณ์ผู้เล่น (User Experience Evaluation):} ประเมินผ่านแบบสอบถามหลังการเล่นกับกลุ่มเป้าหมาย เพื่อวัดระดับความรู้สึกดื่มด่ำ (Immersion) รวมถึงความมั่นใจที่เพิ่มขึ้นและความวิตกกังวลที่ลดลงในการใช้ภาษาอังกฤษ
\end{itemize}